\section{Related Literature}
\label{sec:literature}

The paper sits at the conjunction of four literatures.
Enforcement-design work asks how investigative resources should
be sequenced when monitoring is costly and the layer at which
adjudication happens is more expensive than the layer at which
suspicion can be ranked. Bid-coordination tests detect collusion
off bid-function exchangeability conditional on a
researcher-supplied partition between suspected colluders and a
competitive fringe. Bid-distribution screens detect collusion
off within-tender moments of submitted bids. And cover-bidding
studies identify the loser side of collusive arrangements
retrospectively, after a cartel has been adjudicated. We
position the contribution in the first tradition: an award-layer
triage stage that sequences which firms and environments enter
the bid-layer forensic interrogation. The other three
literatures supply the empirical objects the architecture
operates on---a partition the bid-coordination tests can be
applied to, a forensic battery the bid-distribution screens
implement, and a type the cover-bidding literature has
characterised \emph{ex post} that we identify \emph{ex ante}
on award-record data.

\citet{becker1968crime,stigler1970optimum} establish the
optimal-deterrence baseline: when monitoring is imperfect, the
sanction-detection trade-off has a structural solution that
depends on the informational technology available to the
enforcer. \citet{baker2002case} treats the optimal sequencing of
investigative stages under information asymmetry between agency
and target: the screening rule that sorts candidates into
investigation should economize on data the agency observes at
low marginal cost. \citet{harrington2008detecting} surveys the
detection technologies available to antitrust authorities and
the conditions under which each is informationally feasible.
Two implications travel to our setting. First, when bid-level
observability is absent at the operational layer, the
appropriate screening object is whatever statistic the
collusive arrangement generates that survives the
data-coarsening step. Second, screening and forensic stages are
sequenced, not parallel: the screening stage produces candidate
environments for the forensic stage, which then deploys the
more data-intensive bid-distribution moment battery
\citep{chassang2022robust}. The architectural separation
between an award-layer screening stage and a bid-layer forensic
stage that the construct instantiates is the empirical
implication of this sequencing logic. A growing empirical
literature documents how procurement institutions shape what
enforcement can observe
\citep{coviello2017tenure,decarolis2020bureaucratic,decarolis2020rules};
we treat these institutional features as variation in the
monitoring regime that the construct's signal is read against
(\S\ref{sec:results_oversight}).

\citet{bajari2003deciding} establish the canonical framework for
detecting collusion through exchangeability of bid functions and
conditional independence of residuals, conditional on a
researcher-supplied partition between suspected colluders and a
competitive fringe. The papers that build on Bajari--Ye inherit
the constraint:
\citet{porter1993detection,porter1999ohio} detect bid rotation
conditional on assumed cartel membership;
\citet{baldwin1997bidder}, \citet{clark2021collusion}, and
\citet{schurter2020identification} use post-prosecution records
to identify colluders. The partition is constructed
\emph{ex post} or by assumption in every case; the construct
supplies a candidate \emph{ex ante} firm-level flag computable
from award records that bid-coordination tests can take as input.
\citet{conley2016detecting} is the closest precedent in data
parsimony: operating on Italian highway-construction auctions
where investigators had identified a focal cartel but lacked a
complete membership list, they infer bidder coalitions from
co-bidding network structure. Three differences delimit the
present contribution: their object is \emph{group structure}
(which firms belong to the same coalition given a focal
investigation), while ours is \emph{firm-level flags} on a
screening basis where no investigation has been opened; their
exercise operates on a post-investigation pool, ours on the
\emph{ex ante} universe of always-loser participants; their
analysis recovers within-cartel structure on the assumption that
a cartel exists, ours asks whether the participation footprint
suggests a cartel-affected environment in the first place. The
two methods are sequential, not substitutive.

\citet{imhof2019detecting,imhof2018screening,huber2019machine}
and \citet{wallimann2023machine} train machine-learning
classifiers on tender-level bid statistics---within-tender
coefficient of variation, kurtosis, skewness, normalized
spread, second-lowest distance, and related
moments---reaching AUC above $0.85$ against ground-truth European
cartel cases; \citet{abrantesmetz2006a} pioneer the variance
screen and \citet{chassang2022robust} provide moment-based
theoretical foundations. We re-implement the seven-feature
Imhof--Wallimann pipeline against our cobidder ground truth in
\S\ref{sec:forensic}; isolated single-moment benchmarks (the
within-tender coefficient of variation alone) are strawmen of
the literature's actual practice and we do not use them as the
comparison object. The two layers stand in an asymmetric
observability relation: bid-distribution screens require
per-bidder bid amounts at the analytical-warehouse layer the
screening pipeline runs on, but in most enforcement environments
those amounts are forensic-recoverable through administrative
request rather than routinely operationally queryable
(\S\ref{sec:institutional}); the construct runs on the
operational layer and integrates with bid-distribution screens
at the forensic stage where microdata are recoverable.

\citet{porter1999ohio,pesendorfer2000study,asker2010leniency} and
\citet{marshall2012economics} document cover bidders in
prosecuted cartels through case records, and
\citet{athey2011comparing} compare cover-bidding incentives
across auction formats theoretically; in every empirical case the
cover bidders are identified \emph{ex post}, through leniency
testimony, post-conviction records, or judicial discovery, after
the cartel has been adjudicated. We are not aware of a published
prospective identification of cover-bidder candidates from public
award-record data alone. The bid-level signature implied by the
type varies across this literature:
\citet{porter1993detection,bajari2003deciding} read the
constraint $b_C > b^*$ loosely, with cover bidders placing
\emph{visibly} uncompetitive bids that manufacture the appearance
of competition; \citet{marshall2012economics} and
\citet{asker2010leniency} read it tighter, with cover bids
remaining plausible enough to disguise coordination and
cross-bid dispersion arising from rotation through cover roles.
The bid-level extension of the within-stratum bridge in
\S\ref{sec:cade_bridge} discriminates between the two readings
and is consistent with the second; the formal distinction is
recorded in Online Appendix~A.
